\chapter{The Cross-Platform Contract}
\label{ch:cross-platform-contract}

CNA's two-level taxonomy names four platform families and, inside \texttt{Desktop}, four host
OS outcomes. That is richer than a flat platform list, and support is not a Boolean attached
to any leaf. Linux x86\_64 is the only self-hosting first-class target at the audited revision.
Other paths are cross-compiled and exercised through Wine, a browser, or an Android emulator,
or have not executed at all. The architectural contract is nevertheless coherent: SDL absorbs
common host services, while platform-only code is selected as CMake targets rather than
scattered behind preprocessor branches.

\section{A platform claim is an evidence vector}

The useful progression is:

\begin{enumerate}
  \item a source path exists;
  \item configuration admits it;
  \item it cross-compiles;
  \item it links a real artifact;
  \item the artifact runs under emulation or translation;
  \item it runs on real hardware; and
  \item its output is verified by an appropriate oracle.
\end{enumerate}

Later steps do not follow automatically from earlier ones. A renderer may pass the platform
gate and then fail because a sibling checkout, SDK, prebuilt archive, or 32-bit ABI is missing.
A PE executable may link and still never engage the intended DXVK DLL. A browser test may run
and prove DOM state without proving compositor pixels. This book therefore uses qualified verbs
rather than collapsing the vector into ``supported.''

At this revision no CNA binary is evidenced on genuine Windows, macOS, iOS, or Android hardware.
Windows results run under Wine/DXVK/vkd3d on a real Linux GPU; Web results run in headless
Chromium; Android evidence comes from an AVD and survives mainly as narrative; macOS has no
post-adaptation native build; iOS/tvOS is explicitly unvalidated.

\section{Two compile-time taxonomy levels}

\texttt{CNA::Platform} distinguishes \texttt{Desktop}, \texttt{Android}, \texttt{iOS}, and
\texttt{Web}. A constexpr selector uses Emscripten, Android, and Apple target macros.
\texttt{CNA::DesktopOS} then distinguishes \texttt{Windows}, \texttt{Linux},
\texttt{MacOSX}, and \texttt{Other}, and rejects use on a non-desktop platform.

These enums are real API vocabulary, not a platform service layer. The runtime platform module
has no implementation body, and there is no family of \texttt{cna-platform-*} targets or
pluggable \texttt{IDevicesBackend} objects. Consumers use the taxonomy for decisions such as
sensor selection and orientation support; they do not route all filesystem, window, and device
operations through one universal CNA interface.

\section{CMake exclusion does most of the branching}

Only 56 of 1{,}195 production files under \texttt{modules/} contain an OS platform
preprocessor directive. Those files contain 159 directives, 104 of them in renderer native-
handle or SDK glue. Seven whole modules --- storage, media, audio, math, gamer services,
devices-ext, and graphics-ext --- contain none; graphics has two guarded files among 377.

The low count is structural. Renderer selection rejects impossible platform/identity pairs
before the module is added. A Glide source can include \texttt{windows.h} without an internal
guard because its target is never compiled off 32-bit Windows. There are 109
\texttt{CNA\_RENDERER\_*} conditionals in production, reflecting the true implementation axis,
and 104 platform conditions in CMake that decide which sources exist in the build.

\begin{lstlisting}
# Architectural pattern, simplified:
if(CNA_GRAPHICS_RENDERER STREQUAL "GLIDE")
    # Earlier platform gates have already required 32-bit Windows.
    add_subdirectory(modules/renderers/glide)
endif()
\end{lstlisting}

This is not the same as saying CNA has no platform-specific code. It means excluded code is not
parsed or linked, so it does not need a defensive \texttt{\#ifdef} around every Win32 or Android
token.

\section{SDL is the common platform layer}

Production code contains roughly 5{,}983 SDL/Mixer/Image/TTF symbol occurrences across 125
files --- about 38 SDL API references for every OS preprocessor directive. SDL supplies windows,
events, timers, paths, power information, URL launching, user folders, GL procedure lookup, and
the entry-point bridge. Subsystems link SDL privately instead of routing through a central CNA
platform target.

The abstraction is intentionally not perfect. Launch parameters use Win32 APIs on Windows and
\texttt{/proc/self/cmdline} on Linux/Android, returning empty elsewhere. Graphics-adapter PCI
IDs come from Linux \texttt{/sys/class/drm} and remain zero on other hosts. Glide and GDI use
raw Win32 dynamic loading, while GL families use \texttt{SDL\_GL\_GetProcAddress}. These narrow
exceptions are easier to audit precisely than a nominal universal interface that leaks the same
differences through virtual methods.

\section{Concurrency is mostly supplied by dependencies}

CNA constructs exactly one production \texttt{std::thread}: the Android sensor worker. Audio
registers callbacks into SDL3\_mixer's thread; content loading is synchronous; networking is
single-threaded by contract; input delivery is main-thread-only. Internal synchronization uses
standard \texttt{mutex}, \texttt{atomic}, and \texttt{condition\_variable}, not SDL thread
primitives.

This distinction matters on platforms whose host controls the event loop. A dependency-owned
callback thread is still concurrent with CNA state even though CNA did not create it. Conversely,
the existence of threading primitives in sharp-runtime does not make CNA content or networking
asynchronous.

\section{Renderer admissibility is only configuration truth}

Platform gates admit 26 renderer identities on Linux, 39 on 64-bit Windows, 40 on 32-bit
Windows, 26 on Darwin, 24 on Android/iOS, and 24 on Emscripten. Those numbers mean only that
\texttt{RendererSelection.cmake} does not issue a platform \texttt{FATAL\_ERROR}. They do not
mean dependencies resolve or artifacts build.

Examples expose the gap. OpenGL4 still needs a system OpenGL package; OPENGLES1 needs
\texttt{libGLESv1\_CM}; five GL profiles need easy-gl and meta-gl sibling checkouts;
FREEDIRECT needs free-direct; SKIA needs a prepared artifact; DIRECTX8 needs D8VK import
libraries. No repository artifact enumerates a smaller, fully built set for every platform.

The correct workflow is therefore two-stage: use the gate to reject structurally impossible
pairs, then record dependency, link, execution, engagement, and pixel evidence for the actual
pair tested. Chapters~\ref{ch:windows-cross-compiling} and~\ref{ch:wine-engagement} apply this
discipline to the best-developed non-native path.

\section{The honest platform baseline}

Linux x86\_64 has native execution and real pixel goldens, though its principal workflows are
not uniformly healthy. MinGW Windows has cross-build and translated execution evidence, not
native Windows proof. Web has five renderer identities at four distinct evidence levels.
Android has emulator history without committed APK, logcat, or screenshots, and graphics was
never attempted. Metal's portable policy tests can pass on Linux without compiling Objective-C++
or touching a Metal device. iOS/tvOS remains a hard, explicit non-claim.

This baseline is not pessimism. It prevents a portable helper test, configure-time allowance,
or emulator narrative from being promoted into hardware validation, while preserving the real
engineering each path has already earned.
